\documentclass[11pt, a4paper]{common/document_template}
\usepackage{common}

\begin{document}

\begin{titlepage}
  \begin{center}
	\Huge{Уральский турнир юных математиков}
  \end{center}
\end{titlepage}

\chapter{Номер 46. Год 2015}


\chapter{Номер 49. Год 2017}

\section{Тур 2}

\begin{exercise}
Клетки квадрата $n \times n$ окрашены в чёрный и белый цвета в шахматном порядке. За один ход можно выделить любой квадрат $2 \times 2$ и поменять в нём цвет всех 4 клеток на противоположный. При каких $n$ за несколько таких ходов можно все клетки сделать одноцветными?
\end{exercise}

\begin{Solution}
Легко указать последовательность для перекраски квадрата $4 \times 4$:
\begin{align*}
\begin{pmatrix}
0 & 1 & 0 & 1 \\
1 & 0 & 1 & 0 \\
0 & 1 & 0 & 1 \\
1 & 0 & 1 & 0
\end{pmatrix} \implies
\begin{pmatrix}
1 & 0 & 0 & 1 \\
0 & 1 & 1 & 0 \\
0 & 1 & 0 & 1 \\
1 & 0 & 1 & 0
\end{pmatrix} \implies
\begin{pmatrix}
1 & 0 & 0 & 1 \\
0 & 1 & 1 & 0 \\
0 & 1 & 1 & 0 \\
1 & 0 & 0 & 1
\end{pmatrix} \implies
\begin{pmatrix}
1 & 1 & 1 & 1 \\
0 & 0 & 0 & 0 \\
0 & 1 & 1 & 0 \\
1 & 0 & 0 & 1
\end{pmatrix} \implies
\begin{pmatrix}
1 & 1 & 1 & 1 \\
0 & 0 & 0 & 0 \\
0 & 0 & 0 & 0 \\
1 & 1 & 1 & 1
\end{pmatrix} \implies
\begin{pmatrix}
1 & 1 & 1 & 1 \\
1 & 1 & 1 & 1 \\
1 & 1 & 1 & 1 \\
1 & 1 & 1 & 1
\end{pmatrix}
\end{align*}
Из этого следует, что можно перекрасить любой квадрат, где $n$ - четное. При нечетном $n$ при любой перекраске четность разности между количествами белых и черных клеток не меняется и первоначально нечетная. Значит в этом случае клетки нельзя сделать одноцветными.
\end{Solution}

\chapter{Номер 51. Год 2018}

\section{Тур 1}

\begin{exercise}
У клетчатого квадрата $11 \times 11$ отметили центры всех клеток и 81 из них покрасили в красный цвет. Докажите, что есть такие три красные точки, что одна из них является серединой отрезка с концами в двух других.
\end{exercise}

\begin{Solution}
   
\end{Solution}

\begin{exercise}
Какое наименьшее количество коней могут побить все клетки доски $7 \times 7$, включая клетки, на которых стоят кони? (Сам конь клетку, на которой стоит, не бьёт.)
\end{exercise}

\chapter{Номер 52. Год 2018}

\section{Тур 1}

\begin{exercise}
В классе организовали несколько кружков. В каждый кружок ходит ровно 9 учеников, а каждый ученик - ровно в 2 кружка. Оказалось, что для любых двух кружков есть ровно один ученик, который ходит в оба эти кружка. Сколько детей посещают кружки?
\end{exercise}

\chapter{Номер 54. Год 2019}

\section{Тур 1}

\begin{exercise}
Докажите, что для любого натурального числа $n$ все натуральные числа от 1 до $n$ можно расставить в таком порядке $a_{1}, a_{2}, \dots, a_{n}$, что все числа $a_{1} + 1, a_{2} + 2, \dots, a_{n} + n$ будут натуральными степенями двойки.
\end{exercise}


\chapter{Номер 55. Год 2020}

\section{Тур 1}

\begin{exercise}
11 школьников написали тест. Учитель проверил работы и заметил, что для любых двух вопросов теста найдутся хотя бы 6 школьников, каждый из которых ответил правильно ровно на один из этих двух вопросов. Докажите, что в тесте было не более 12 вопросов.
\end{exercise}

\begin{Solution}
Обозначим количество вопросов как $n$. Рассмотрим количество троек $(q_{i}, q_{j}, s_{k})$, где $q_{i}$ и $q_{j}$ это два вопроса, а $s_{k}$ - это школьник, который ответил правильно ровно на один из этих двух вопросов. С одной стороны это количество равно $6 \binom{n}{2}$. С другой оно равно
$$
\sum_{k = 1}^{11} p_{k} (n - p_{k})
$$
где $p_{k}$ - количество вопросов на которые $k$-ый школьник ответил верно. Несложно заметить, что эта сумма достигает максимума если $n = 2p_{k}$. Значит
$$
6\binom{n}{2} \leqslant \frac{11 n^{2}}{4} \implies 12(n - 1) \leqslant 11n 
$$
Откуда $n \leqslant 12$.
\end{Solution}

\begin{exercise}
В стране 100 сенаторов, у каждого из которых ровно 4 помощника. Сенаторы и их помощники состоят в комитетах. Комитет может состоять либо из 5 сенаторов, либо из 4 сенаторов и 4 помощников, либо из 2 сенаторов
и 12 помощников. Каждый сенатор состоит в пяти комитетах, а каждый помощник - в трёх. Сколько всего комитетов в стране? (Разумеется, один
и тот же человек не может быть помощником сразу у двух сенаторов.)
\end{exercise}

\begin{Solution}
Легко видеть, что всего $100 \cdot 4 = 400$ помощников. Теперь пусть $X$, $Y$, $Z$ - количество комитетов первого, второго и третьего типа. Тогда подсчитаем пары губернатор-комитет и помощник-комитет двумя способами.
$$
\begin{cases}
5X + 4Y + 2Z = 5 \cdot 100 \\
4Y + 12Z = 3 \cdot 400
\end{cases}
$$
Умножив первое уравнение на 4 и сложив со вторым получим $20(X + Y + Z) = 3200$. Откуда $X + Y + Z = 160$.
\end{Solution}

\chapter{Номер 57. Год 2021}

\section{Тур 1}

\begin{exercise}
В некоторые $k$ клеток доски $8 \times 8$ расставлены различные числа таким образом, что для каждой клетки с числом $a$ среди чисел в клетках, имеющих ровно одну общую точку с данной, не более одного числа, большего $a$. Найдите наибольшее возможное значение $k$.
\end{exercise}

\begin{exercise}
Есть 6 гирь, веса которых 10, 10, 11, 11, 11, 12 килограмм. На каждой из них наклейка, на которой указан вес этой гири. Эксперт знает, что наклейки правильные, но его шеф сомневается - он считает, что его сын мог ради шутки оторвать некоторые наклейки и переклеить их по-другому. Как эксперту всего за 2 взвешивания на
двухчашечных весах убедить шефа, что наклейки всё-таки правильные?
\end{exercise}

\begin{exercise}
Учитель объявил 99 умным детям, что загадал 99 последовательных натуральных чисел, одно из которых делится на 1000. Он сообщил каждому по секрету одно из этих 99 чисел (все сообщенные числа различны). Затем учитель раздал всем по бумажке и попросил каждого из детей написать, может ли он определить сумму 99 загаданных чисел. Каждый написал "Нет". Учитель собрал бумажки и объявил
детям, что все написали "Нет", после чего снова раздал бумажки и продолжал так делать, пока кто-то впервые не напишет "Да". Может ли так случиться, что все 99 детей напишут "Да" одновременно?
\end{exercise}


\chapter{Номер 58. Год 2022}

\section{Тур 4}

\begin{exercise}
Игорь отметил 768 точек на окружности. Оказалось, что длины дуг между
всеми парами соседних точек равны $1, 2, 3, \dots, 768$ в некотором порядке. Какое наибольшее количество пар диаметрально противоположных точек могло получиться у Игоря?
\end{exercise}

\begin{exercise}
Вова выбрал бесконечное множество $S$ целых чисел. Он называет целое число чистым, если оно представляется в виде суммы нескольких различных чисел из $S$ единственным образом, и грязным в противном случае (то есть если представляется несколькими разными способами или вообще не представляется). Докажите, что грязных чисел либо не существует, либо бесконечно много.
\end{exercise}


\chapter{Номер 59. Год 2022}

\section{Тур 1}

\begin{exercise}
Даны 100 различных действительных чисел. Докажите, что их можно расставить в клетках таблицы $10 \times 10$ так, чтобы для любых двух чисел, стоящих в соседних по стороне ячейках, их разность была не равна 1.
\end{exercise}

\section{Тур 2}

\begin{exercise}
Назовём полоской клетчатый прямоугольник, хотя бы одна из сторон которого
равна 1. Мы хотим разбить клетки квадрата $99 \times 99$ на $m$ непересекающихся полосок так, чтобы любые две из этих m полосок имели общую границу длины не более 1. Полоски могут быть разных длин. Определите наименьшее $m$, для которого это возможно.
\end{exercise}

\begin{exercise}
На доске $21 \times 42$ на некоторые клетки поставлены камни. На доску ставится робот, который умеет ходить в каждом из четырёх направлений. Однако менять направление робот может, только уперевшись в границу доски или в камень. Какое наименьшее количество камней надо поставить на доску, чтобы робот, стартовав с некоторой клетки в определенном направлении, мог обойти все клетки доски, не занятые камнями? В одной и той же клетке робот может побывать несколько раз.
\end{exercise}

\section{Тур 3}

\begin{exercise}
Каждая клетка полоски $1 \times 2022$ покрашена в синий или красный цвет. Назовем блоком последовательность идущих подряд одноцветных клеток, слева и справа от которой находятся или концы полоски, или клетки другого цвета. Краб, находящийся в некоторой клетке полоски, может сделать ход по следующему правилу: если он сидит в блоке размера $m$ и этот блок синий, то он перемещается на $m$ клеток вправо, а если блок красный, то на $m$ клеток влево. Если такой ход выводит за пределы полоски, то ход не совершается и движение заканчивается. При каком наименьшем $k$ можно закрасить $k$ клеток полоски синим цветом, а остальные - красным так, чтобы краб за несколько ходов мог добраться из самой левой клетки
в самую правую?
\end{exercise}


\section{Тур 4}

\begin{exercise}
Множества $A$ и $B$ содержат по 750 натуральных чисел из диапазона от 1 до 2022 включительно. Может ли так быть, что ни одно из чисел вида $a + b$, где $a \in A$, $b \in B$, не является точным квадратом? 
\end{exercise}


\chapter{Номер 60. Год 2023}

\section{Тур 2}

\begin{exercise}
Докажите, что существуют различные натуральные $a$ и $b$, большие 1000000, для которых $a^{b} + 1$ делится на $b^{a} + 1$.
\end{exercise}

\begin{exercise}
Зайнаб изготовила ожерелье из 1734 бусин, на них в некотором порядке написаны числа $290, 291, \dots, 2023$. Докажите, что в этом ожерелье найдутся три идущие подряд бусины, числа на которых являются длинами сторон треугольника.
\end{exercise}

\section{Тур 3}

\begin{exercise}
Квадрат $99 \times 99$ разрезали на квадраты со сторонами $1, 2, \dots, 9$, все девять размеров использованы. Будем называть ряд клеток (горизонталь или вертикаль) нечётным, если он пересекает нечётное число квадратиков. Найдите наименьшее возможное количество нечётных рядов.
\end{exercise}

\section{Тур 4}

\begin{exercise}
Последовательность $a_{n}$ задана рекуррентным соотношением $a_{0} = 1$, $a_{1} = x + 1$, $a_{n+2} = xa_{n+1} - a_{n}$, где $x$ - натуральное число. Докажите, что найдется бесконечно много $x$, для которых в последовательности нет простых чисел.
\end{exercise}

\chapter{Номер 61. Год 2023}

\section{Командная олимпиада}

\begin{exercise}
В каждой клетке прямоугольной таблицы $1000 \times 1000$ стоит рыцарь или лжец. Рыцарь всегда говорит правду, а лжец всегда лжёт. Каждый заявил, что в клетках, соседних с его клеткой по стороне, стоит поровну лжецов и рыцарей. Может ли на доске быть ровно 2023 рыцаря?
\end{exercise}

\section{Тур 2}

\begin{exercise}
Оля расставила по кругу 200 пустых мисок. Раз в секунду она кладет кусочек
корма в одну из мисок по своему выбору, после чего кошка Кайя сдвигает лапкой один кусочек (не обязательно тот, который только что добавили!) из любой миски в соседнюю. Кайя хочет, чтобы через 200 секунд в какой-нибудь миске оказалось хотя бы 10 кусочков корма. Может ли она действовать так, чтобы гарантированно добиться желаемого?
\end{exercise}

\section{Тур 3}

\begin{exercise}
На столе лежат 54 кучки с $1, 2, \dots, 53, 54$ камнями. Разрешается выбрать любую кучку, подсчитать количество $k$ имеющихся в ней камней и убрать по $k$ камней из всех кучек, в которых камней не меньше $k$. В результате нескольких таких операций на столе осталась одна кучка. Сколько камней в ней может быть?
\end{exercise}

\section{Тур 4}

\begin{exercise}
Дано натуральное $n$. Докажите, что можно закрасить некоторые клетки бесконечной клетчатой плоскости в зелёный цвет так, чтобы любой прямоугольник из $n$ клеток содержал нечётное число зелёных клеток.
\end{exercise}


\chapter{Номер 62. Год 2024}

\section{Тур 1}

\begin{exercise}
Докажите, что если натуральное число делится на 88, но не делится на 121, то все его делители можно разбить на два множества с одинаковой суммой.
\end{exercise}

\begin{exercise}
Копы и роббер бегают по конечному связному графу. Сначала копы занимают некоторые вершины. После них роббер подбирает себе начальную вершину. Ходят по очереди: все копы, потом снова роббер, потом снова все копы и так далее. Ход каждого персонажа - это либо пас (остаться на месте), либо переместиться в соседнюю вершину. Все всех видят, действия копов скоординированы. Цель копов - поймать
роббера, это происходит если в какой-либо момент роббер и любой из копов окажется в одной вершине. Докажите, что существует граф, на котором роббер может неограниченно долго ускользать от 2024 копов.
\end{exercise}


\begin{exercise}
В первой строке треугольной таблицы из 1002 строк написаны в порядке возрастания числа $0, 1, 2, \dots, 1001$. Каждая строка, начиная со второй, строится по такому правилу: под промежутком между каждыми двумя соседними числами предыдущей строки записывается их сумма. Докажите, что единственное число в последней строке делится на 1001.
\end{exercise}

\begin{Solution}
Несложно доказать по индукции, что если в первой строке таблицы стоят числа $a_{0}, a_{1}, a_{2}, \dots, a_{n}$, то в последней строке будет стоять число
$$
S = \sum_{k = 0}^{n} \binom{n}{k} a_{k}
$$
Тогда в нашем случае
$$
S = \sum_{k = 0}^{n} k \binom{n}{k} = n 2^{n - 1}
$$
\end{Solution}

\section{Тур 2}

\begin{exercise}
Числа $2024, 2025, 2026, \dots, 2124$ записаны подряд в некотором порядке, образуя многозначное число. Докажите, что оно не простое.
\end{exercise}

\begin{Solution}
Это число можно представить в виде
$$
S = \sum_{i = 0}^{100} a_{i} 10000^{i} = \sum_{i = 0}^{100} a_{i} \mod 101
$$
Так как $10000 = 1 \mod 101$. Тогда 
$$
S = \sum_{i = 0}^{100} a_{i} \mod 101 =  \sum_{i = 0}^{100} i \mod 101 = 5050 \mod 101 = 0
$$
\end{Solution}

\begin{exercise}
Два брата, решив заработать денег, пошли к сельскому психологу. Тот запер их в две отдельные комнаты и стал играть с ними в следующую игру. Время от времени психолог заходит в комнату к одному из братьев и спрашивает его, какой по счету этот заход (психолог считает заходы к обоим братьям вместе), на что вопрошаемый дает два ответа. Если хотя бы один из ответов верный, то психолог зачисляет один рубль на счет угадывающего, но ему об этом не сообщает. Братья
могут заранее договориться об образе действий. Могут ли они гарантированно заработать хотя бы 100 рублей?
\end{exercise}

\section{Тур 3}


\section{Тур 4}

\begin{exercise}
Докажите, что существует бесконечно много натуральных $n$ таких, для которых числа $1, 2, 3, \dots, n$ можно раскрасить в два цвета так, чтобы произведение попарных сумм чисел первого цвета было равно произведению попарных сумм чисел второго цвета.
\end{exercise}

\begin{Solution}
Для $n = 4$ мы можем покрасить 1 и 4 в первый цвет, а 2 и 3 во второй. Теперь пусть существует искомая покраска чисел от 1 до $n = 2^{m}$. Тогда покрасим все числа от $n + 1$ до $2n$ так, что $n + i$, где $1 \leqslant i \leqslant n$, покрашено в цвет, противоположный цвету $i$. Например для $n = 4$ получаем одноцветные множества $S_{1}$ и $S_{2}$:
$$
S_{1} = (1, 4, 6, 7) \quad S_{2} = (2, 3, 5, 8)
$$
Докажем, что полученная покраска чисел от 1 до $2n$ удовлетворяет условиям задачи. Очевидно, количество чисел первого цвета равно количеству чисел второго. Пусть $a_{1}, a_{2}, \dots, a_{N}$ будет последовательностью чисел, покрашеных в первый цвет, а $b_{1}, b_{2}, \dots, b_{N}$ - во второй.
\end{Solution}

\begin{exercise}
В полном графе на 100 вершинах ребра покрашены в несколько цветов так, что не имеется разноцветных треугольников. Все ребра, выходящие из вершины $A$, покрашены в разные цвета. Докажите, что в графе найдется такая вершина $B$, что все ребра, выходящие из этой вершины, покрашены (ровно) в 66 цветов.
\end{exercise}


\chapter{Номер 63. Год 2024}

\section{Тур 1}

\begin{exercise}
Отмечены 100 узлов на левой стороне квадрата $100 \times 100$ и 100 узлов его на нижней стороне, причем левый нижний узел не отмечен. Саша синим карандашом нарисовал $n$ четырехугольников с отмеченными вершинами так, что никакие два не имеют общих точек. То же самое сделал красным карандашом сделал Игорь. При каком наименьшем $n$ обязательно найдутся красный и синий четырехугольник, не имеющие общих точек?
\end{exercise}

\begin{exercise}
В начале на доске записаны числа 1 и 2. За одну операцию можно выбрать любое положительное число $x$ и заменить числа $a$, $b$ на доске на числа $a + \frac{x}{b}$ и $b + \frac{x}{a}$. Могут ли
после несколько операций на доске оказаться числа 2 и 3?
\end{exercise}

\begin{Solution}
Рассмотрим отношение $a$ к $b$. После замены оно будет равно
$$
\frac{a + \frac{x}{b}}{b + \frac{x}{a}} = \frac{a}{b}
$$
Значит это отношение является инвариантом. Но исходное отношение равно $\frac{1}{2}$, а конечное $\frac{2}{3}$. Значит 2 и 3 на доске оказаться не могут.
\end{Solution}

\begin{exercise}
Петя посчитал на калькуляторе, что первая цифра $2^{999}$ равна 5. Докажите, что вторая цифра числа $5^{999}$ больше 5.
\end{exercise}

\begin{exercise}
Известно, что $ab + bc + ca = 1$. Докажите, что
$$
\frac{a - b}{c^{2} + 1} + \frac{b - c}{a^{2} + 1} + \frac{c - a}{b^{2} + 1} = 0
$$
\end{exercise}

\begin{Solution}
Преобразуем первый член тождества
\begin{align*}
\frac{a - b}{c^{2} + 1} = \frac{a - b}{c^{2} + ab + bc + ca} = \frac{a - b}{(c + a)(c + b)} = \frac{1}{c + b} - \frac{1}{c + a}
\end{align*}
Аналогично
$$
\frac{b - c}{a^2 + 1} = \frac{1}{a + c} - \frac{1}{a + b}, \quad 
\frac{c - a}{b^2 + 1} = \frac{1}{b + a} - \frac{1}{b + c}
$$
Отсюда следует искомое.
\end{Solution}

\begin{exercise}
Двузначное число называется специальным, если оно является началом десятичной записи чисел $4^{n}$ и $5^{n}$ для некоторого натурального $n$. Известно, что существуют ровно два специальных числа. Чему они равны?
\end{exercise}

\begin{exercise}
В некоторой компании людей для любых двух человек имеется ровно один третий, кого они оба уважают. (Уважение не обязательно взаимно.) Докажите, что каждого человека уважает столько же людей, скольких уважает он сам.
\end{exercise}

\begin{exercise}
Саша и Паша по очереди закрашивают клетки квадрата $12 \times 12$ в красный и синий цвет. Закрашивать дважды одну и ту же клетку запрещено. Начинает Саша. За один ход он красит одну клетку, а Паша - сразу три. Саша выигрывает, если по окончании игры найдется диагональ, в которой количества красных и синих клеток отличаются
хотя бы на 2, иначе выигрывает Паша. У кого из игроков есть выигрышная стратегия?
\end{exercise}

\begin{exercise}
Можно ли среди чисел $1, 2, 3, \dots, 1000$ выбрать 10 так, чтобы сумма никаких двух выбранных чисел не делилась на сумму никаких двух других выбранных чисел? Например, числа 1, 2, 4 вместе брать нельзя, так как 2 + 4 делится на 1 + 2.
\end{exercise}

\section{Тур 2}

\begin{exercise}
Докажите, что при каждом натуральном $n \geqslant 5$
$$
S(1)S(3)\dots S(2n-1) \geqslant S(2)S(4)\dots S(2n)
$$
где $S(k)$ - сумма цифр числа $k$.
\end{exercise}

\begin{exercise}
У натурального числа $n \geqslant 2$ выписали все делители:
$$
1 = d_{1} < d_{2} < \dots < d_{k} = n
$$
Могут ли в этом ряду два соседних числа $d_{s}$ и $d_{s+1}$ оказаться точными квадратами?
\end{exercise}

\begin{Solution}
Пусть $p_{1}, p_{2}, \dots, p_{k}$ будут простыми делителями $n$. Допустим нашлись два таких соседних делителя $d_{s} = a^{2}$ и $d_{s+1} = b^{2}$, где
$$
a = p_{1}^{\alpha_{1}} p_{2}^{\alpha_{2}} \dots p_{k}^{\alpha_{k}}, \quad b = p_{1}^{\beta_{1}} p_{2}^{\beta_{2}} \dots p_{k}^{\beta_{k}}
$$
Тогда выберем делитель $c = ab$
$$
c = p_{1}^{\alpha_{1} + \beta_{1}} p_{2}^{\alpha_{2} + \beta_{2}} \dots p_{k}^{\alpha_{k} + \beta_{k}}  
$$
Очевидно, что $d_{s} < ab < d_{s+1}$. И что $c$ является делителем
$$
p_{1}^{2 \max(\alpha_{1}, \beta_{1})} p_{2}^{2 \max(\alpha_{2}, \beta_{2})} \dots p_{k}^{2 \max(\alpha_{k}, \beta_{k})}
$$
Значит $c | n$ и $c$ лежит между $d_{s}$ и $d_{s+1}$.
\end{Solution}

\begin{exercise}
Клетки таблицы $3 \times 100$ покрашены в черный и белый цвет. Множество клеток называется связным, если от любой его клетки до любой другой можно добраться, переходя несколько раз в соседнюю по стороне клетку, принадлежащую этому множеству. Докажите, что существует связное множество, у которого количества чёрных и белых клеток отличаются не меньше чем на 100.
\end{exercise}

\begin{exercise}
В 329 комнатах находится 329 человек. Все они разного роста. Раз в минуту из каждой комнаты уходит самый высокий человек и переходит в другую комнату (по своему усмотрению), а все остальные остаются на месте. Правда ли, что с помощью этих действий из любого размещения людей по комнатам может получиться любое другое?
\end{exercise}

\section{Тур 3}

\begin{exercise}
На плоскости нарисован 222-угольник и три различные прямые. Каждая из этих прямых проходит ровно через $n$ вершин 222-угольника. Найдите наибольшее возможное значение $n$.
\end{exercise}

\begin{exercise}
Несколько школьников писали олимпиаду из 6 задач, за каждую можно получить от 0 до 7 баллов. Оказалось, что невозможно выбрать двух школьников и 5 задач так, что по каждой из этих 5 задач эти школьники набрали поровну баллов. Какое наибольшее количество школьников могли писать олимпиаду?
\end{exercise}

\begin{exercise}
Докажите, что существует возрастающая арифметическая прогрессия длины 2024 с натуральных членами, \\ отличающимися друг от друга только перестановкой цифр.
\end{exercise}

\begin{Solution}
Циклические числа.
\end{Solution}

\section{Тур 4}

\begin{exercise}
Дано натуральное число $n$. Клетки доски $4n \times 4n$ покрашены в 4 цвета. Множество из 4 клеток назовем пёстрым, если эти клетки покрашены в 4 разных цвета и находятся на пересечении двух строк и двух столбцов таблицы. Какое наибольшее число пёстрых множеств может быть в этой таблице?
\end{exercise}

\begin{exercise}
Муравей ползает по поверхности куба по замкнутому маршруту. Оказалось,
что он всегда ползёт параллельно какому-то ребру куба, поворачивает только под прямым углом и, попадая на ребро куба, поворачивает так, чтобы ползти по следующей грани, но не по ребру. Рассмотрим прямолинейные отрезки его маршрута (от поворота до следующего поворота). Для каждой грани посчитали количество таких отрезков. Могло ли так случиться, что эти шесть количеств образуют шесть последовательных натуральных чисел в каком-то порядке?
\end{exercise}

\chapter{Номер 64. Год 2025}

\section{Тур 1}

\begin{exercise}
При каких вещественных $x$, $0 < x < 180$, равносторонний треугольник можно разбить на конечное число треугольников, каждый из которых имеет угол $x$ градусов?
\end{exercise}

\section{Тур 2}

\begin{exercise}
Может ли в выпуклом 21-угольнике быть не менее тридцати трёх диагоналей длины 1?
\end{exercise}


\end{document}
